% !Mode:: "TeX:UTF-8"
% 任务书中的信息
%% 原始资料及设计要求
\assignReq
{1. AutoChip原论文}
{2. VerilogEval、RTLLM评估基准}
{3. Icarus Verilog开源EDA工具}
{4. 大语言模型API（Claude、GPT系列）}
{5. 北航本科毕业设计相关规范}
%% 工作内容
\assignWork
{1. 复现AutoChip风格RTL代码生成与反馈修复系统}
{2. 接入VerilogEval和RTLLM双benchmark}
{3. 在三个不同能力等级的LLM上验证反馈循环有效性}
{4. 通过消融实验分离反馈信息与多采样的贡献}
{5. 探索反馈粒度、多轮对话和提示词策略的影响}
{6. 建立实验产物管理与审计机制}
%% 参考文献
\assignRef
{[1] Thakur S, et al. AutoChip: Automating HDL Generation, 2023}
{[2] Liu M, et al. VerilogEval: Evaluating LLMs for Verilog, ICCAD 2023}
{[3] Lu Y, et al. RTLLM: Open-Source Benchmark for RTL Generation, 2024}
{[4] Liu S, et al. RTLCoder: LLM-Assisted RTL Code Generation, 2024}
{[5] Tsai Y, et al. RTLFixer: Fixing RTL Syntax Errors with LLMs, 2024}
{[6] Yao Z, et al. HDLDebugger: Streamlining HDL Debugging, 2025}
{[7] Wei J, et al. Chain-of-Thought Prompting, NeurIPS 2022}
{[8] Brown T, et al. Language Models are Few-Shot Learners, NeurIPS 2020}
